\chapter{Preliminaries}

\begin{definition}[Überschneidung]
	Zwei Kanten gelten als Schneidet wenn sie sich in der Ebene kreuzen.
	Im Folgendem werden mit der Funktoin \(I: ( \mathbb{Z}^2, \mathbb{Z}^2) \longrightarrow \{\mathbb{Z}^2\}\) alle Kanten angegeben, die die übergebene Kante schneiden.
\end{definition}
\begin{definition}[Planarer Graph]
	Ein Planarer Graph ist ein Graph ohne überschneiden Kanten.
	Das bedeutet man kann alle Punkte und Kanten darstellen ohne das sich diese überschneiden.
\end{definition}
\begin{definition}[Triangulation]
	Eine Triangulation einer Punktemenge \(P \in \mathbb{Z}^2\) ist ein Maximaler Planarer Graph.
	Das bedeutet man kann keine weiteren Kanten hinzufügen ohne das sich diese Schneiden.
\end{definition}
\begin{theorem}
	TODOGehört Sowas hier hin? Soll ich sowas nennen


	Alle Flächen außer der Äußeren Fläche einer Triangulation bestehen aus Dreiecken.
	Der Beweis wird von Mark Berg et.al. geliefert\cite[p.~185]{Delaunay_Triangulation}.
\end{theorem}
\begin{definition}[Grad]
	Der Grad gibt die Anzahl der anliegenden Kanten an einem Knoten an.
	Im folgenden wird er durch die Funktion \(\degree: \mathbb{Z}^2 \longrightarrow \mathbb{Z}\) angegeben.
\end{definition}
\begin{definition}[Gradbeschränkte Triangulation]
	Die Gradbeschränkte Triangulation hat die gleichen Vorraussetzungen wie eine Triangulation.
	Zusätzlich gibt es noch eine Funktion \(\degrees: \mathbb{Z}^2 \longrightarrow \mathbb{Z}\).
	Diese gibt den Grad an welcher ein Knoten haben muss.
\end{definition}
\begin{definition}[Flipp]
	Eine Operation bei der eine Kante entfernt und eine andere Hinzugefügt wird.
	Hierbei dürfen nach der Operation die Triangulations Eigenschaften nicht verletzt sein.
	Ein Flip wird durch die Funktion \(\flip:(\mathbb{Z}^2,\mathbb{Z}^2) \longrightarrow (\mathbb{Z}^2,\mathbb{Z}^2)\) definiert.

\end{definition}
\begin{definition}[Diagonale Flipp]
	Ein Diagonaler Flipp hat die gleichen Eigenschaft wie ein Flipp.
	Zusätzlich wird für die Kante \((b,c)\) die leeren Dreicke \(\triangle abc \) und  \(\triangle bcd \) gesucht und \((a,d)\) zurückgeben.
	Er wird durch die Funktion \(\flipDio:(\mathbb{Z}^2,\mathbb{Z}^2) \longrightarrow (\mathbb{Z}^2,\mathbb{Z}^2) definiert\).
\end{definition}
